% !TITLE 赋得古原草送别
% !AUTHOR 白居易
% !DYNASTY 唐
% !GENRE 五言律诗
% !ORDER 44

\begin{poem}
离离\textsuperscript{1}原上草，一岁一枯荣\textsuperscript{2}。
野火烧不尽，春风吹又生。
远芳侵古道\textsuperscript{3}，晴翠接荒城\textsuperscript{4}。
又送王孙去\textsuperscript{5}，萋萋满别情\textsuperscript{6}。
\end{poem}

\begin{annotations}
\notetext{1}{离离：草木繁盛茂密的样子。}
\notetext{2}{一岁一枯荣：每年枯萎一次，繁盛一次。枯荣二字写尽了生命的轮回。}
\notetext{3}{远芳侵古道：远处的草香蔓延到古道上。侵，侵占、蔓延。荒废的古道被芳草占领。}
\notetext{4}{晴翠接荒城：晴日里的翠绿色连接到荒城边。接，连接、延伸到。}
\notetext{5}{王孙：贵族子弟的泛称，这里指远行的友人。典出《楚辞》"王孙游兮不归，春草生兮萋萋"。}
\notetext{6}{萋萋：草茂盛的样子，也谐音"凄凄"（伤感）。满别情，草有多茂盛，离别之情就有多浓。}
\end{annotations}

\begin{appreciation}
这是一首命题诗。唐代科举考试中，考官给定题目，考生现场作诗，称为"赋得"。白居易写这首诗时年仅十六岁，却已显露出不凡的气象。

离离原上草，一岁一枯荣——离离是草木茂盛的样子。原野上的青草，每年秋天枯萎，春天又繁盛起来。枯荣二字，写尽了草木的生命轮回，也暗喻人生的聚散离合。一个十六岁的少年，能从中读出这样的况味，眼光已超越同龄人。

野火烧不尽，春风吹又生——这是全诗最著名的一联，也是中国人最熟悉的诗句之一。野火燎原，看似毁灭了一切，但根还在土里；春风一吹，又冒出嫩芽。这两句的力量在于对比：野火是暴烈的、短暂的，春风是温和的、持久的；毁灭是外在的，生命是内在的。白居易没有直接说"生命顽强"，而是用一幅画面让人自己体会。这种写法，比任何议论都更有说服力。

远芳侵古道，晴翠接荒城——远处的草香蔓延到古道上，晴日里的翠色连接到荒城边。侵和接两个字，写出了青草不可阻挡的生长力。古道和荒城是人世的遗迹，芳草是自然的新生，两者相对，有一种时光流转、物是人非的感慨。

又送王孙去，萋萋满别情——王孙原指贵族子弟，这里泛指友人。萋萋是草盛的样子，也谐音"凄凄"。满别情三字，将前六句对草的描写全部收束到送别之情上。草越茂盛，别情越浓；草越绵延，思念越长。

这首诗的妙处，在于以草写人，以景写情。表面上句句写草，实际上句句写送别。白居易后来以"文章合为时而著，歌诗合为事而作"闻名，主张诗歌要反映现实。但这首少年时代的作品告诉我们，他的才华首先在于对自然和情感的敏锐捕捉。野火烧不尽，春风吹又生——这不仅是草的精神，也是所有在困境中不放弃希望的人的精神。
\end{appreciation}
